\section{Component Ablation Study}
\label{sec:ablation}

To quantify the individual marginal utility of each architectural component within HydroEdge-Net, we conducted a systematic ablation study isolating the biophysical feature engine, categorical context encodings, and model structural constraints.

\subsection{Ablation Configurations and Empirical Results}

We evaluated four systematic variants under identical training splits and evaluation protocols:
\begin{enumerate}
    \item \textbf{Variant 1 (Full Proposed HydroEdge-Net):} Incorporating the complete feature vector $\mathbf{x} \in \mathbb{R}^8$ (including Tetens VPD, SMD, $\Psi_{\text{TMS}}$, soil texture, and phenology encodings) with constrained depth ($d \le 5, N_{\text{trees}}=45$).
    \item \textbf{Variant 2 (w/o Biophysical Features):} Ingesting only raw physical sensor telemetry ($T, \text{RH}, \theta$, crop type, soil type, phenology) without derived thermodynamic formulations (omitting VPD, SMD, and $\Psi_{\text{TMS}}$).
    \item \textbf{Variant 3 (w/o Soil and Phenology Encodings):} Ingesting only ambient and edaphic variables ($T, \text{RH}, \theta, \text{VPD}, \text{SMD}, \Psi_{\text{TMS}}$) while discarding categorical soil texture and phenological stage encodings.
    \item \textbf{Variant 4 (Unconstrained Heavy GBDT):} Retaining the full feature set but removing all microcontroller structural bounds ($d = 12, N_{\text{trees}} = 100$), representing uncompressed cloud-grade ensembles.
\end{enumerate}

Table~\ref{tab:ablation_study} catalogs the quantitative performance across these ablation variants on the held-out test split ($N = 3,283$).

\begin{table*}[!t]
\caption{Systematic Component Ablation Study on Held-Out Test Split ($N = 3,283$)}
\label{tab:ablation_study}
\centering
\renewcommand{\arraystretch}{1.25}
\resizebox{0.85\textwidth}{!}{
\begin{tabular}{lcccc}
\toprule
\textbf{Architectural Variant} & \textbf{Accuracy} & \textbf{Weighted F1-Score} & \textbf{Deficit Recall} & \textbf{Model Flash Footprint} \\
\midrule
\textbf{1. Full Proposed HydroEdge-Net} & \textbf{95.71\%} & \textbf{0.9518} & \textbf{98.96\%} & \textbf{16.8~KB} \\
\textbf{2. w/o Biophysical Features} (Raw sensor inputs only) & 95.86\% & 0.9541 & 99.04\% & 28.0~KB \\
\textbf{3. w/o Soil \& Phenology Encodings} (Ambient data only) & 89.67\% & 0.8755 & 97.11\% & 12.4~KB \\
\textbf{4. Unconstrained Heavy GBDT} ($d=12, N_{\text{trees}}=100$) & 99.97\% & 0.9997 & 99.92\% & 185.0~KB \\
\bottomrule
\end{tabular}
}
\end{table*}

\subsection{Key Ablation Insights and Discussion}

The ablation results provide decisive empirical validation of our engineering design:
\begin{enumerate}
    \item \textbf{Decisive Impact of Soil Texture and Phenology (Variant 1 vs. 3):} Discarding categorical soil texture and phenological development stages precipitates a dramatic performance collapse: overall accuracy plummets by \textbf{6.04\%} (from 95.71\% to 89.67\%), and weighted F1 drops by \textbf{0.0763} (from 0.9518 to 0.8755). This drop confirms that identical soil moisture readings ($\theta = 22\%$) carry fundamentally disparate physiological implications depending on soil matric suction (e.g., Sandy soil vs. Heavy Clay) and crop phenological vulnerability (e.g., shallow-rooted Seedling vs. deep-rooted Maturation). Phenological awareness is mathematically indispensable for reliable edge diagnosis.
    \item \textbf{Biophysical Features as Silicon Compressors (Variant 1 vs. 2):} Training the model on raw sensor telemetry alone (Variant 2) yields a comparable diagnostic accuracy (95.86\% vs. 95.71\%), but inflates the compiled Flash memory footprint by \textbf{66.67\%} (from 16.8~KB to 28.0~KB). When denied explicit thermodynamic equations, the decision tree must generate multiple levels of nested piecewise-constant split nodes to approximate the exponential curvature of the Tetens vapor pressure curve and the multiplicative coupling of $\Psi_{\text{TMS}}$. Providing closed-form biophysical features directly in the preprocessing step relieves the tree of learning complex non-linear thermodynamics, allowing HydroEdge-Net to achieve equivalent accuracy with a \textbf{40.0\% reduction in Flash memory}.
    \item \textbf{Radical Silicon Compression vs. Heavy Cloud Ensembles (Variant 1 vs. 4):} The unconstrained heavy ensemble (Variant 4) achieves near-perfect test accuracy (99.97\%), but demands 185.0~KB of Flash storage. HydroEdge-Net achieves near-parity performance (95.71\% accuracy and 98.96\% deficit recall) while slashing Flash memory consumption by \textbf{90.92\%} (16.8~KB vs. 185.0~KB). This confirms that HydroEdge-Net captures over 95\% of optimal diagnostic capacity while fitting comfortably into the most restrictive microcontroller tiers.
\end{enumerate}
